% Generated draft: requires mathematical and source review.
\documentclass[11pt]{article}
\usepackage[margin=25mm]{geometry}
\usepackage{fontspec}
\setmainfont{Latin Modern Roman}
\usepackage{amsmath,amssymb,mathtools,unicode-math}
\setmathfont{Latin Modern Math}
\usepackage{enumitem,needspace,xurl,hyperref}
\setlength{\parindent}{0pt}
\setlength{\parskip}{5pt}
\setlength{\emergencystretch}{3em}
\newcommand{\N}{\mathbb{N}}
\newcommand{\Z}{\mathbb{Z}}
\newcommand{\Q}{\mathbb{Q}}
\newcommand{\R}{\mathbb{R}}
\newcommand{\C}{\mathbb{C}}
\newcommand{\F}{\mathbb{F}}
\newcommand{\A}{\mathbb{A}}
\newcommand{\PP}{\mathbb{P}}
\newcommand{\E}{\mathbb{E}}
\newcommand{\eps}{\varepsilon}
\newcommand{\dd}{\,\mathrm{d}}
\begin{document}
\section*{UnsolvedMath ID: 30001450}
\textbf{Source problem code:} catalogue entry
\subsection{Definitions and mathematical statement}
Let \(\mathfrak g\) be a symmetrizable Kac--Moody algebra, let \(\Lambda\) be a dominant integral weight, and let \(q\) be an indeterminate. Write \(U_q(\mathfrak g)\) for the corresponding quantized enveloping algebra over \(\Q(q)\), and interpret \(V(\Lambda)\) as its irreducible integrable highest-weight module of highest weight \(\Lambda\).

The record denotes by \(K^\Lambda(\Gamma)\) the Grothendieck group of graded finite-dimensional modules over the cyclotomic quiver Hecke algebra. To make the usual interpretation explicit, after specifying a coefficient field and the defining algebra parameters, let \(R^\Lambda(\beta)\) denote the cyclotomic quiver Hecke algebra at \(\beta\in Q_+\), where \(Q_+\) is the positive root monoid of \(\mathfrak g\), and put
\[
K^\Lambda(\Gamma)=\bigoplus_{\beta\in Q_+}K_0\bigl(R^\Lambda(\beta)\text{-}\mathrm{gmod}_{\mathrm{fd}}\bigr).
\]
Here \(\mathrm{gmod}_{\mathrm{fd}}\) denotes the category of finite-dimensional graded modules with degree-preserving morphisms, and \(K_0\) imposes the relations \([M]=[L]+[N]\) for short exact sequences \(0\to L\to M\to N\to0\). With the convention that \(q[M]=[M\langle1\rangle]\), where \(\langle1\rangle\) is a grading shift, this is a module over \(\Z[q,q^{-1}]\). These conventions explain the displayed scalar extension; the record itself does not specify them or define \(\Gamma\).

\textbf{Supplied question, with unresolved qualifications.} For every such \(\mathfrak g\) and \(\Lambda\), does there exist an isomorphism of \(U_q(\mathfrak g)\)-modules
\[
\Phi:\Q(q)\otimes_{\Z[q,q^{-1}]}K^\Lambda(\Gamma)\xrightarrow{\ \sim\ }V(\Lambda),
\]
under which appropriately normalized classes of graded simple modules correspond to the dual canonical basis of \(V(\Lambda)\)?

This is an intended formulation rather than a fully specified conjecture: the coefficient field, defining parameters, meaning of \(\Gamma\), quantum-group action, and grading normalization of simple modules are absent. The phrase ``dual canonical basis'' also requires a precise basis construction and any hypotheses under which the asserted correspondence is intended. Without a normalization, all grading shifts of a simple module cannot individually correspond to distinct elements of the stated basis.

The supplied dated assessment describes the underlying extraction as an incomplete progress paragraph about the Khovanov--Lauda conjecture and retains it as open under an unclear-fragment rule. Thus the displayed question must not be treated as an independently verified open problem, nor may the supplied reference to a later proof be taken to resolve every clause and generality of this formulation.
\subsection{Short English statement}
The intended question asks whether graded finite-dimensional modules over cyclotomic quiver Hecke algebras recover the highest-weight representation \(V(\Lambda)\) through their Grothendieck group, with normalized simple-module classes giving its dual canonical basis. Essential algebra, action, and basis conventions are missing. The supplied open label reflects dated curation of an incomplete fragment, not an independently verified mathematical status.
\subsection{Sources}
\begin{itemize}
\item Local UnsolvedMath registry, permanent ID \(30001450\), source problem code \texttt{catalogue entry}, title ``Cyclotomic Quiver-Hecke Categorification of Highest-Weight Modules,'' in the Oberwolfach Reports Open Problems collection. The supplied proposed year is 2010; no individual proposer is identified.
\item Supplied report citation: \emph{Combinatorial Representation Theory} (2010), Oberwolfach Reports, pp.~799--882, DOI \texttt{10.4171/owr/2010/15}. Original OWR report: \url{https://doi.org/10.4171/owr/2010/15}. The supplied background attributes the material to mathematical workshop discussions and connects it to the Khovanov--Lauda conjecture.
\item Supplied literature source labeled ``later proof of the identifiable cyclotomic theorem'': \url{https://arxiv.org/abs/1102.4677}. No authors, title, or precise theorem scope for this reference are supplied in the record.
\item Literature assessment checked on \texttt{2026-08-21}: supplied dated curation, not independently verified here. It labels the record open because the extracted progress paragraph is incomplete, while also listing a later proof of an identifiable cyclotomic theorem. The supplied curation expressly requests independent verification of time-sensitive status and cited claims.
\item The registry labels the statement \texttt{corrected\_verified} and reports checking it against the official report. This is a supplied verification claim; no source was independently visited or checked in preparing this draft.
\end{itemize}
\end{document}
